\documentclass{article}

\usepackage{amsmath, amssymb, bm}
\usepackage[a4paper, margin=1in]{geometry}

\title{Drone control system}
\author{}
\date{}

\begin{document}
	
	\maketitle
	\section{Introduction}
	This document describes the basic control system of our drone.
	
	\section{Conventions}
	We will make some conventions here. Our quadcopter drone motors are number as shown below (TODO insert figure).
	
	Other conventions:
	\begin{itemize}
		\item Motor 1,3 spin counterclockwise
		\item Motor 2,4 spin clockwise
	\end{itemize}
	
	\section{Mixer}
	The mixer is in charge of taking the torque vector $T$ and required vertical force component $F_z$ and producing the required thrust vectors $F_{t1}$, $F_{t2}$, $F_{t3}$ and $F_{t4}$.
	
	In a subsequent step we can translate these thrust vectors to rpm setpoints because we know the relation between thrust and speed is:
	\[
	F_t = k_{t}\cdot\omega^2
	\]
	Where $k_{t}$ is a property of the propeller.
	
	Determining $T_x$ and $T_y$ is straight forward:
	\begin{align*}
	T_x = F_{t1}\cdot l_1 + F_{t4}\cdot l_1 - F_{t2}\cdot l_1 - F_{t3}\cdot l_1 \\
	T_y = F_{t3}\cdot l_2 + F_{t4}\cdot l_2 - F_{t1}\cdot l_2 - F_{t2}\cdot l_2 
	\end{align*}
	
	Too determine $T_z$ we need too consider that when the propeller spins due to rotation resistance a pure reaction torque is generated opposing the spinning direction. The reaction torque has following relation to the rpm:
	\[
		T_d = k_d\cdot\omega^2
	\]
	We can define:
	\[
		c = \frac{T_d}{F_t} = \frac{k_d}{k_t}
	\]
	We can then write the torque $T_z$ in relation to the propeller forces:
	\[
	T_z = c\cdot(F_{t1} +  F_{t3} - F_{t2} - F_{t4})
	\]
	We know the total thrust $F$ is simply the addition of all the propeller forces:
	\[
	F = F_{t1} + F_{t2} + F_{t3} + F_{t4}
	\]
	However since our higher level controllers are outputting vertical force we need the relation between thrust and vertical force:
	\[
	F_z = F\cdot\cos(\phi)\cdot\cos(\theta)
	\] 
	Where $\phi$ is the roll angle and $\theta$ is the pitch angle.
	
	Now we can arrange all equations into a matrix:
	\[
	\begin{bmatrix}
		F \\
		T_x \\
		T_y \\
		T_z
	\end{bmatrix} = 
	\begin{bmatrix}
		1 & 1 & 1 & 1 \\
		l_1 & -l_1 & -l_1 & l_1 \\
		-l_2 & -l_2 & l_2 & l_2 \\
		c & -c & c & -c
	\end{bmatrix}\cdot
	\begin{bmatrix}
		F_{t1} \\
		F_{t2} \\
		F_{t3} \\
		F_{t4}
	\end{bmatrix}
	\]
	
	Or written differently:
	\[
	u = A \cdot f
	\]
	We now have a simple linear system of equations and we can find the required propeller forces by:
	\[
	f = A^{-1}\cdot u 
	\]
\end{document}